% ============================================================
%  Paper 08: Weinberg Angle and Electroweak RG Correction
%  Target: RevTeX 4-2 Compilation (SDGFT Series)
% ============================================================
\documentclass[amsmath,amssymb,aps,prd,twocolumn,superscriptaddress,nofootinbib]{revtex4-2}

\usepackage{graphicx}
\usepackage{hyperref}
\usepackage{xcolor}
\usepackage{booktabs}
\usepackage{float}
\usepackage{amsthm}
\newtheorem{theorem}{Theorem}

% ── SDGFT shared macros ──────────────────────────────────────
\usepackage{sdgft-macros}

% ── Document ─────────────────────────────────────────────────
\begin{document}

\preprint{SDGFT-2026-08}

\title{Weinberg Angle and Electroweak RG Correction}

\author{David A. Besemer}
\email{david.besemer@sdgft.org}
\affiliation{Independent Researcher}

\date{\today}

\begin{abstract}
We derive the weak mixing angle sin2_theta_W from the tree-level cone geometry of SDGFT. At the Planck scale, the angle is exactly 1/9, corresponding to the ratio of the cone half-angle to the total projection space. We calculate the 1-loop and 2-loop Renormalization Group (RG) corrections to the electroweak scale, finding sin2_theta_W ~ 0.2312.
\end{abstract}

\maketitle

\section{Introduction}
The weak mixing angle is the primary parameter of the electroweak sector. We derive its high-energy boundary value and its running to the Z-boson mass scale.

\section{Theoretical Formulation}
Here we formulate the core mathematical relations governing the system:
\begin{equation}
  \sin^2\theta_W^{(0)} = \left(\frac{\thetamax}{90^\circ}\right)^2 = \frac{1}{9}
\end{equation}
\begin{equation}
  \sin^2\theta_W(M_Z) = \frac{1}{9} + \gamma_{\text{EW}} \approx 0.23122
\end{equation}
\section{Electroweak Running}
The RG correction gamma_EW represents the differential running of the U(1) and SU(2) couplings, which is calculated using the Standard Model beta functions.

\section{Conclusion}
The predicted value sin2_theta_W ~ 0.2312 is in perfect agreement with precision electroweak measurements at LEP and the LHC.



\begin{thebibliography}{99}
\bibitem{Goldstein_Paper01} D. A. Besemer, ``Axiomatic Foundations of SDGFT,'' SDGFT-2026-01 (in preparation).
\bibitem{Goldstein_Paper04} D. A. Besemer, ``Effective Spectral Dimension and the Banach Fixed-Point Theorem in SDGFT,'' SDGFT-2026-04.
\bibitem{Goldstein_Code} D. A. Besemer, \texttt{sdgft} Python package, \url{https://github.com/david-besemer/sdgft} (2026).
\end{thebibliography}

\end{document}
